(i) Fact \claimref{odd-extensions.imported-minimal-completion-degenerate-algebra-rule} with $(p,q,r) = (1,3,1)$ gives $(2,4)$. An isometry of $\widetilde W$ fixing $W$ pointwise sends $\varphi$ to a vector $x$ with $Q(x, w) = Q(\varphi, w)$ for all $w \in W$ and $Q(x,x) = 0$; the five linear conditions have rank five in six unknowns, so $x = \varphi + cI$, and $Q(x,x) = 2c\,Q(I,\varphi) = 2c$ forces $c = 0$. (ii) $Q(aI + b\varphi, aI + b\varphi) = 2ab\,Q(I,\varphi) = 2ab$; the vector is orthogonal to the Lorentz frame because $I$ and $\varphi$ are. (iii) The plane $\mathrm{span}(I,\varphi)$ is the orthogonal complement of the Lorentz frame, so an isometry fixing the frame restricts to an isometry of the plane, and the isometry group of a hyperbolic plane is $O(1,1)$: the boosts $(a,b) \mapsto (ta, b/t)$, $t > 0$, the swap $(a,b) \mapsto (b,a)$, and their products with $-1$. The boosts are transitive on each branch of each hyperbola $2ab = \text{const} \neq 0$ and on each of the four null rays; the element $-1$ exchanges the two branches of a hyperbola, and the swap exchanges the two null axes; so the orbits are the level sets of the norm. (iv) Isometries preserve the norm; the non-isomorphism is Theorem \claimref{odd-extensions.twistor-algebra}(iv). The script verifies the signature, the rank-five linear system and its unique null solution, the norm formula on instances of both signs and zero, the boost and swap invariances, and the centres.
